\documentclass[10pt,a4paper]{amsart}
\usepackage{amsmath,amssymb,amsthm,graphicx}
\usepackage{verbatim}
\newenvironment{reference}{\comment}{\endcomment}
%\newenvironment{reference}{}{}
\newenvironment{slogan}{\comment}{\endcomment}
\newenvironment{history}{\comment}{\endcomment}

% For commutative diagrams we use Xy-pic
\usepackage[all]{xy}

% We use 2cell for 2-commutative diagrams.
\xyoption{2cell}
% Theorem environments.
%
\theoremstyle{plain}
\newtheorem{theorem}[subsection]{Theorem}
\newtheorem{proposition}[subsection]{Proposition}
\newtheorem{lemma}[subsection]{Lemma}

\theoremstyle{definition}
\newtheorem{definition}[subsection]{Definition}
\newtheorem{example}[subsection]{Example}
\newtheorem{exercise}[subsection]{Exercise}
\newtheorem{situation}[subsection]{Situation}

\theoremstyle{remark}
\newtheorem{remark}[subsection]{Remark}
\newtheorem{remarks}[subsection]{Remarks}

\numberwithin{equation}{subsection}

% Macros
%
\def\lim{\mathop{\mathrm{lim}}\nolimits}
\def\colim{\mathop{\mathrm{colim}}\nolimits}
\def\Spec{\mathop{\mathrm{Spec}}}
\def\Hom{\mathop{\mathrm{Hom}}\nolimits}
\def\Ext{\mathop{\mathrm{Ext}}\nolimits}
\def\SheafHom{\mathop{\mathcal{H}\!\mathit{om}}\nolimits}
\def\SheafExt{\mathop{\mathcal{E}\!\mathit{xt}}\nolimits}
\def\Sch{\mathit{Sch}}
\def\Mor{\mathop{\mathrm{Mor}}\nolimits}
\def\Ob{\mathop{\mathrm{Ob}}\nolimits}
\def\Sh{\mathop{\mathit{Sh}}\nolimits}
\def\NL{\mathop{N\!L}\nolimits}
\def\CH{\mathop{\mathrm{CH}}\nolimits}
\def\proetale{{pro\text{-}\acute{e}tale}}
\def\etale{{\acute{e}tale}}
\def\QCoh{\mathit{QCoh}}
\def\Ker{\mathop{\mathrm{Ker}}}
\def\Im{\mathop{\mathrm{Im}}}
\def\Coker{\mathop{\mathrm{Coker}}}
\def\Coim{\mathop{\mathrm{Coim}}}

% Boxtimes
%
\DeclareMathSymbol{\boxtimes}{\mathbin}{AMSa}{"02}

%
% Macros for moduli stacks/spaces
%
\def\QCohstack{\mathcal{QC}\!\mathit{oh}}
\def\Cohstack{\mathcal{C}\!\mathit{oh}}
\def\Spacesstack{\mathcal{S}\!\mathit{paces}}
\def\Quotfunctor{\mathrm{Quot}}
\def\Hilbfunctor{\mathrm{Hilb}}
\def\Curvesstack{\mathcal{C}\!\mathit{urves}}
\def\Polarizedstack{\mathcal{P}\!\mathit{olarized}}
\def\Complexesstack{\mathcal{C}\!\mathit{omplexes}}
% \Pic is the operator that assigns to X its picard group, usage \Pic(X)
% \Picardstack_{X/B} denotes the Picard stack of X over B
% \Picardfunctor_{X/B} denotes the Picard functor of X over B
\def\Pic{\mathop{\mathrm{Pic}}\nolimits}
\def\Picardstack{\mathcal{P}\!\mathit{ic}}
\def\Picardfunctor{\mathrm{Pic}}
\def\Deformationcategory{\mathcal{D}\!\mathit{ef}}
\usepackage{cleveref}
\crefname{originalreference}{original reference}{original references}
\Crefname{originalreference}{Original reference}{Original references}
\makeatletter
\newlabel{lemma-normalization-completion}{{0BG9}{source}{}{sourcepointer-1}{}}
\newlabel{lemma-normalization-completion@cref}{{[originalreference][0][]0BG9}{[1][0][]source}{}{}{}}
\makeatother
\title{Descending sequences of normalized blowups along the completion}
\author{The Stacks Project Authors}
\date{stacks-project:resolve.tex:local-snapshot-20261012; GNU FDL 1.2 or later}
\begin{document}
\maketitle
\noindent\textit{Editorial note.} All delivered mathematical text is the author's own, copied from the pinned Stacks Project chapter file resolve.tex; no statement or proof step was rewritten, reordered or removed.\par\medskip
\noindent\textit{Editorial note.} The counted proof is the author's complete argument: the induction matching closed points through the isomorphism of closed fibres, and the descent of the normalization step along the completion.\par\medskip
\noindent\textit{Editorial note.} The label lemma-normalized-blowup-completion names the counted statement, and the proof environment immediately following it in the source is delivered verbatim as its proof.\par\medskip
\noindent\textit{Editorial note.} The counted statement is carried by the author's own lemma environment in the chapter source.\par\medskip
\noindent\textit{Editorial note.} Pre-open gate: python operations/agent-authoring/tools/precard\_gate.py --source resolve.tex --label lemma-normalized-blowup-completion --start 3235 --end 3262 --local-id 28111 returned RESULT=OK (L28111 free, occupancy FREE, public interval clear). Pre-assembly tag self-check: canonical tag resolved and not present in reference/public-items.jsonl.\par\medskip
\noindent\textit{Editorial note.} The statement and proof are copied from the pinned chapter file; the cited inputs are kept as external pointers with their canonical Stacks tags.\par\medskip
\noindent\textit{Allowed standard prerequisites (editorial).}\begin{itemize}
\item Compatibility of normalization with completion for an excellent local domain. Same chapter, the lemma on normalization and completion. Transfers the normalization step of the sequence over the completion back to the original ring, giving finiteness of the normalization.
\item Base change of blowing up in a closed point and identification of the special fibres. Divisors, the compatibility of blowing up with flat base change, and the isomorphism of the closed fibres of the members of the sequence with those of their base changes. Shows that the blowing up over the completion is the base change of the blowing up of the corresponding closed point over the original ring.
\end{itemize}\medskip
\ifdefined\hypertarget\hypertarget{sourcepointer-1}{}\fi
\noindent\textit{Source reference.} Original reference 0BG9: sources/stacks-resolve/pinned-original-source.tex, source lines 3190--3190. Prerequisite pointer invoked at proof line(s) 3259; it names a standard fact of the same chapter that the counted argument uses, and every step of the counted proof built on it is retained verbatim in the delivered body.\par\medskip
\begin{lemma}
\label{lemma-normalized-blowup-completion}
Let $(A, \mathfrak m, \kappa)$ be a Noetherian local domain whose completion
$A^\wedge$ is normal. Then given any sequence
$$
Y_n \to Y_{n - 1} \to \ldots \to Y_1 \to \Spec(A^\wedge)
$$
of normalized blowups, there exists a sequence of (proper) normalized blowups
$$
X_n \to X_{n - 1} \to \ldots \to X_1 \to \Spec(A)
$$
whose base change to $A^\wedge$ recovers the given sequence.
\end{lemma}

\begin{proof}
Given the sequence $Y_n \to \ldots \to Y_1 \to Y_0 = \Spec(A^\wedge)$ we
inductively construct $X_n \to \ldots \to X_1 \to X_0 = \Spec(A)$.
The base case is $i = 0$. Given $X_i$ whose base change is $Y_i$,
let $Y'_i \to Y_i$ be the blowing up in the closed point $y_i \in Y_i$
such that $Y_{i + 1}$ is the normalization of $Y_i$.
Since the closed fibres of $Y_i$ and $X_i$ are isomorphic, the point
$y_i$ corresponds to a closed point $x_i$ on the special fibre of $X_i$.
Let $X'_i \to X_i$ be the blowup of $X_i$ in $x_i$. Then the base change
of $X'_i$ to $\Spec(A^\wedge)$ is isomorphic to $Y'_i$. 
By Lemma \ref{lemma-normalization-completion}
the normalization $X_{i + 1} \to X'_i$ is finite and its base change
to $\Spec(A^\wedge)$ is isomorphic to $Y_{i + 1}$.
\end{proof}

\end{document}
